\section{Техническое задание}

\begin{enumerate}
    \item Основание для выполнения работы (учебный план)
    \item Исполнитель и соисполнители работы (студент)
    \item Цели, задачи и исходные данные
        \begin{enumerate}[label=\alph*.]
            \item Целью работы является разработка программного комплекса на языке Haskell для статистического анализа растровых шрифтов в формате \verb|BDF|, позвояющего выявлять глифы с аномальными метриками и низкой читаемотью, чтобы предотвратить ошибки отображения данных на утсройствах c ограниченными ресурсами и низким разрешением дисплеев, например, как счётчики, тепминалы оплаты и бытовая электроника.
            \item В процессе выполнения работы требуется решить следующие задачи.
            \begin{enumerate}
                \item Изучить спецификацию формата растровых шрифтов BDF.
                \item Разработать парсер файлов формата BDF на языке Haskell.
                \item Реализовать математические модели для оценки метрик глифов.
                \item Разработать алгоритм поиска визуально похожих глифов.
                \item Реализовать интерфейс для взаимодействия с пользователем.
                \item Обеспечить логирование действий и обработку ошибок ввода.
                \item Покрыть ключевые модули тестами свойств с использованием библиотеки QuickCheck.
                \item Подготовить отчетную документацию.
            \end{enumerate}
            \item Исходными данными для проведения работы являются.
            \begin{enumerate}
                \item Методические указания по выполнению курсовой работы.
                \item Техническая документация формата \verb|BDF|.
                \item Файлы растровых шрифтов в формате \verb|.bdf|.
            \end{enumerate}
        \end{enumerate}
    \item Основное содержание работы
    \begin{enumerate}[label=\alph*.]
        \item Разработка проекта.
        \begin{enumerate}
            \item Изучения спецификации формата \verb|BDF| и на её основе проектирования типов данных для представления шрифта и глифов.
            \item Реализация парсера файла шрифта формата \verb|.bpf|.
            \item Реазилация модулей анализа шрифта и метрик глифов.
            \item Интеграция всех медудей в единый программый комплекс.
            \item Покрытие получившегося программного комплекса тестами.
        \end{enumerate}
    \end{enumerate}
\end{enumerate}
